\section{Server API and Security Semantics}

Each API command is specified as a typed transition with explicit checks. The API assumes authenticated TLS sessions as stated in the project description.

\subsection{Data Types}

We use:
\begin{itemize}
    \item \texttt{Principal}: actor in $\mathcal{A}$,
    \item \texttt{ConfLabel}: conjunction of policies $(o:r)$,
    \item \texttt{IntegLabel}: conjunction of policies $(o \leftarrow w)$,
    \item \texttt{Label}: pair $(C,I)$,
    \item \texttt{DataId}, \texttt{ProgramId}: server-generated identifiers.
\end{itemize}

For each command, response is either success value or \texttt{error(code,msg)}.

\subsection{Command List}

\textbf{1) UPLOAD}$(token, payload, C, I, meta)$: Creates object with label $(C,I)$. Preconditions: valid token, valid labels, caller in $\mathcal{A}$. Response: \texttt{ok(dataId)}. Security: validates label syntax (I1 preserved).

\textbf{2) READ\_DATA}$(token, dataId)$: Returns payload if caller in all reader sets of $C$. Security: read checks prevent unauthorized access (I2 preserved).

\textbf{3) WRITE\_DATA}$(token, dataId, payload')$: Updates payload if caller in all writer sets of $I$ and $I \sqsubseteq I'$. Security: write checks prevent untrusted modifications (I3 preserved).

\textbf{4) UPLOAD\_PROGRAM}$(token, src, signature)$: Parses, type-checks, and flow-checks source code. Stores checked IR with inferred interface. Response: \texttt{ok(programId, interface)} or \texttt{error(type/flow\_violation)}. Security: rejects unsafe code before storage (I4 preserved).

\textbf{5) RUN\_PROGRAM}$(token, programId, inputIds, outC, outI)$: Executes program if (1) caller authorized for inputs, (2) computed $L_{req} = L_1 \sqcup \cdots \sqcup L_n \sqsubseteq (outC, outI)$. Security: ensures output label contains all input information (I4 preserved).

\textbf{6) DECLASSIFY\_AND\_RELEASE}$(token, dataId, transform, params, targetC)$: Applies approved transform if (1) transform whitelisted, (2) params satisfy policy (e.g., $k \geq k_{min}$), (3) for relaxed policies $(o:r) \to (o:r')$, caller has authority from owner $o$. Stores result with label $(targetC, I')$ and audit entry. Security: authority-checked declassification (I5 preserved).
    \item \textbf{Security}. DECLASSIFY\_AND\_RELEASE preserves I5 (declassification safety) because no policy is relaxed without authority, and every relaxation is logged.
\end{itemize}

\textbf{7) AUDIT\_QUERY}$(token, filter)$

Retrieves audit log entries subject to access control.

\begin{itemize}
    \item \textbf{Preconditions}: token valid; $filter$ specifies criteria (e.g., date range, actor, $dataId$).
    \item \textbf{Checks}: caller must be an administrator \emph{or} an auditor authorized by an owner mentioned in the filter. If caller is an owner seeking audit entries about their own data, they are authorized.
    \item \textbf{Effect}: no state change.
    \item \textbf{Response}: \texttt{ok(entries)} where each entry is a tuple $(actor, time, source\_dataId, source\_label, target\_label, transform, params, authority\_proof\_reference)$. Or \texttt{error(access\_denied)} if caller not authorized.
    \item \textbf{Security}. AUDIT\_QUERY supports I5 accountability by allowing data owners to verify that declassifications of their data were performed correctly and only with proper authority.
\end{itemize}

\subsection{Why This API Is Secure by Construction}

The interface intentionally separates concerns, and each separation reinforces security:

\begin{itemize}
    \item \textbf{Data commands (UPLOAD, READ, WRITE)}. These enforce reader/writer policies on stored objects. READ and WRITE are the only ways a user can interact with persistent data, and both are gated by label checks. This ensures that every data access respects the stored confidentiality and integrity labels, regardless of which user is making the request or what application logic might be running.
    
    \item \textbf{Program commands (UPLOAD\_PROGRAM, RUN\_PROGRAM)}. These enforce static flow safety before execution. Programs are checked once at UPLOAD time; once a program is stored, its safety is proven and fixed. RUN\_PROGRAM does not re-check the program (trust cache), but does verify that the input and output labels align with the checked interface. This keeps program execution fast while maintaining guarantees.
    
    \item \textbf{Declassification command (DECLASSIFY\_AND\_RELEASE)}. This centralizes all confidentiality relaxation under explicit authority, making the policy boundary clear. Declassification is never implicit or hidden; every relaxation is logged and auditable. No user program can perform declassification; only the server can execute approved transforms.
    
    \item \textbf{Audit command (AUDIT\_QUERY)}. This provides the evidence trail to verify that all declassifications were authorized and that no illegal flows occurred. Owners can audit their own data's history, and administrators can audit system-wide activity.
\end{itemize}

By separating command concerns (data access, program safety, declassification, audit), we prevent checks from interfering and ensure each property composes.